\section{Summary and Future Prospects}
\label{sec:summary}

Section~\ref{sec:results} compared the Conv1D autoencoder with PCA and showed
both the promise and the current limitations of two-number compression.  This
section summarizes the main conclusions and the next steps of the project.

To isolate compression from observational uncertainty, the comparison uses
100,000 controlled surfaces generated by varying 13 sampled shape parameters.
Each surface contains 3,131 values over redshift and wavenumber.  PCA provides
the linear baseline, while the Conv1D autoencoder forces the same surface
through a two-number bottleneck.  Holding the representation length fixed is
what lets the experiment test the value of nonlinearity rather than the value
of storing more numbers.

On the held-out validation set, the autoencoder with a two-number latent space
recovers $98.09\%$ of the total variance, compared with $96.23\%$ for
two-component PCA.  At the same latent dimension, the nonlinear autoencoder
therefore reconstructs the surfaces more accurately than the linear PCA
baseline.  This suggests that the model family contains structure that is not
captured as efficiently by two linear PCA components.  However, these
measurements come from a preliminary short run.  The autoencoder has not met
the $99.9\%$ recovered-variance target or been evaluated on the test set, and
its errors remain sizeable for difficult surfaces.  The result therefore
supports continuing the nonlinear approach, but not yet replacing the
13-parameter description in a cosmological analysis.

The next step is to test larger latent dimensions, which will show how much
accuracy is gained by relaxing the two-number constraint, and to repeat the
most promising configurations with several random seeds, which will test
whether the gain is reproducible.  Once the architecture and latent dimension
have been chosen, the untouched test set can provide the final generalization
check.  In the longer term, the model family should be connected more directly
to physical intrinsic-alignment calculations and survey requirements, because
surface reconstruction is useful only if its errors are small where they
matter for lensing inference.
